\documentclass{article}

\usepackage[final]{Styles/neurips_2025}
\usepackage[utf8]{inputenc}
\usepackage[T1]{fontenc}
\usepackage{hyperref}
\usepackage{url}
\usepackage{booktabs}
\usepackage{amsfonts}
\usepackage{microtype}
\usepackage{xcolor}
\usepackage{enumitem}

\title{Assignment 3 Report: Agent-Oriented LLM Post-Training for Structured Replanning in EmbodiedSceneAgent}

\author{
Che CHENG \\
Student ID: 225015050 \\
Course: CSC5051: Natural Language Processing
}

\makeatletter
\renewcommand{\@noticestring}{}
\makeatother

\begin{document}
\maketitle

\begin{abstract}
This report presents EmbodiedSceneAgent as an Assignment 3 study on agent-oriented LLM post-training for structured replanning with scene-memory grounding, verifier feedback, and failure-aware recovery. The work focuses on a cognition-layer closed loop rather than end-to-end low-level control. The primary quantitative evidence is a directly comparable base-vs-tuned evaluation on structured planner outputs ($n=73$), where the tuned path improves tool-use accuracy, target-match rate, and strict task proxy while maintaining perfect format compliance. A secondary fixed tiny 3-case qualitative comparison under matched settings shows accepted revised plans of $0/3$ (stable baseline), $3/3$ (Qwen2.5-VL-3B), and $2/3$ (Qwen2.5-VL-7B). The report emphasizes conservative interpretation, explicit evidence hierarchy, and current limitations.
\end{abstract}

\section{Introduction and Research Topic}
This project studies language-conditioned embodied high-level planning, where a model must propose and revise executable subtasks under changing scene state. The downstream task is not generic text generation; it is structured decision-making in a closed cognition loop.

The work is an LLM post-training task because model outputs are constrained to a fixed planner schema, and performance is evaluated by comparing a stable baseline against a tuned model under the same contract. It is agent-oriented because planning is embedded inside local executable components: observation processing, scene memory, verifier checks, and replanning.

The main problem addressed is failure recovery after initial planning errors, especially grounding mismatch and repeated no-effect behavior. The practical target is to make recovery auditable and improve structured decision quality without overstating benchmark claims.

\section{Task Formulation}
\textbf{Input.} Natural-language instruction, current observation-derived scene state, and recent action/history context.

\textbf{Intermediate state.} Structured \texttt{SceneMemory} used as a grounded object-level representation.

\textbf{Planner output schema.} \texttt{PlannerOutput} fields:
\texttt{task}, \texttt{subgoal}, \texttt{target\_object}, \texttt{skill}, \texttt{success\_check}, \texttt{fallback}, with optional \texttt{reasoning} and \texttt{confidence}.

\textbf{Verifier and replanner role.} The verifier compares before/after state to classify failure types; the replanner then performs rule or hybrid rule+LLM revision with audit traces.

\textbf{Success/failure definition.} Success means the episode reaches objective completion under loop policy. Failure means effective progress cannot be achieved before termination.

\section{Experiment Design}
\subsection{Model Selection and Comparison Setup}
Base model path in this report: stable baseline planner path used for the main quantitative comparison and tiny qualitative baseline anchor.

Tuned model path in this report: \texttt{Qwen2.5-VL-3B-Instruct} with minimal LoRA SFT.

The tuning artifact is recorded in \texttt{results/checkpoints/planner\_sft\_3b\_minimal/run\_latest/}, with model/training metadata in \texttt{run\_meta.json} and setup snapshot in \texttt{config.snapshot.yaml}. Reported metadata includes \texttt{model\_id=Qwen/Qwen2.5-VL-3B-Instruct}, \texttt{max\_optimizer\_steps=80}, and \texttt{best\_val\_loss=0.4959}.

\subsection{3090-Only Development Strategy}
The current development path uses a constrained setup suitable for RTX 3090:
\begin{itemize}[leftmargin=*]
  \item Minimal LoRA SFT configuration uses bf16 precision with compact low-rank adaptation (\texttt{r=8}, \texttt{alpha=16}, \texttt{dropout=0.05} on \texttt{q\_proj}/\texttt{v\_proj}).
  \item Short-step, small-batch training is explicitly configured (\texttt{max\_steps=80}, \texttt{per\_device\_batch\_size=1}, \texttt{gradient\_accumulation\_steps=4}).
  \item Real checkpoint, run metadata, and config snapshot are versioned as reproducibility artifacts.
  \item No claim is made that this is a large-scale final training regime.
\end{itemize}

\subsection{Data and Supervision Construction}
Training sample shape (report shorthand):
\begin{itemize}[leftmargin=*]
  \item input = instruction + scene memory + history + failure evidence
  \item target = structured revised plan JSON
\end{itemize}
The repository uses structured planner supervision and failure-aware examples rather than unconstrained QA-style text targets. In practical terms, supervision is aligned to the fixed planner contract (\texttt{task, subgoal, target\_object, skill, success\_check, fallback}), so the training/evaluation objective is schema-consistent decision quality instead of free-form fluency.

\subsection{Structured Output and Prompt Contract}
The planner/replanner path enforces a strict key contract and canonical skill vocabulary. Runtime includes parse, repair, validation, and fallback logic; only schema-valid and semantically accepted revised plans are executed.

\subsection{Evaluation Protocol and Evidence Hierarchy}
\textbf{Primary quantitative evidence:} base-vs-tuned metrics from \texttt{results/eval/planner\_base\_vs\_tuned/metrics.json} with $n=73$.

\textbf{Secondary qualitative evidence:} fixed tiny 3-case side-by-side comparison under matched settings (\texttt{backend=calvin\_debug\_real}, \texttt{batch=grouped\_sequence}, \texttt{episodes=3}, \texttt{seed=42}, \texttt{verifier\_mode=verifier\_plus\_replan}, \texttt{replanner\_mode=hybrid}).

\textbf{Engineering evidence:} config snapshots, run manifests, failure taxonomies, and audit fields from generated artifacts.

\section{Code and System Implementation}
\subsection{Closed-Loop Architecture}
The implemented cognition loop is:
\begin{center}
\texttt{Observation -> Scene Memory -> Planner -> Executor -> Verifier -> Replanner}
\end{center}

\subsection{Tool-Like Components}
\begin{itemize}[leftmargin=*]
  \item \textbf{Scene memory interface}: provides grounded object ids and state fields.
  \item \textbf{Planner contract}: enforces structured output for downstream routing.
  \item \textbf{Verifier}: maps execution outcomes to explicit failure evidence.
  \item \textbf{Hybrid replanner}: combines rule-first repair with optional LLM fallback.
\end{itemize}

\subsection{Parser, Repair, and Validation}
Replanning uses a constrained JSON pipeline: extraction/recovery, key sanitization, contract validation, and parse-error labeling. Invalid outputs fall back to safer rule behavior with explicit audit reasons.

\subsection{Semantic Acceptance Filter}
Before executing an LLM-revised plan, the system rejects schema-valid but semantically inconsistent plans (for example, target absent from scene memory or drawer-goal target mismatch). This separates format correctness from grounding correctness.

\subsection{Repeated-No-Effect Guard}
The loop records consecutive no-effect signatures and can stop repeated ineffective fallback attempts early, producing explicit terminal labels for diagnosis.

\subsection{Reproducibility Artifacts}
Run directories include config snapshots and run manifests. This enables traceable experiment provenance without introducing unverified claims.

\section{Experimental Results}
\subsection{Main Quantitative Comparison}
Table~\ref{tab:main-quant} reports the direct base-vs-tuned comparison from the current repository artifacts.
\begin{table}[t]
\caption{Primary quantitative comparison (base vs tuned, $n=73$). Metrics are structured planner-output proxies, not official CALVIN environment success rates.}
\label{tab:main-quant}
\centering
\begin{tabular}{lcccc}
\toprule
Track & Format compliance & Tool-use accuracy & Target-match rate & Strict task proxy \\
\midrule
Stable baseline (planner base) & 1.000 & 0.082 & 0.329 & 0.068 \\
Tuned mainline (LoRA 3B minimal) & 1.000 & 0.219 & 0.479 & 0.205 \\
\bottomrule
\end{tabular}
\end{table}


This table is the primary quantitative evidence in this report because it is directly comparable under one fixed evaluator and one fixed planner-output contract ($n=73$), without mixing different backends or run protocols. Metric meanings are:
\begin{itemize}[leftmargin=*]
  \item \textbf{Format compliance}: whether planner outputs satisfy the required structured format.
  \item \textbf{Tool-use accuracy}: whether the selected skill aligns with reference targets in the proxy evaluation.
  \item \textbf{Target-match rate}: whether the chosen action-target grounding matches reference targets.
  \item \textbf{Strict task proxy}: conjunction-style proxy requiring format correctness, skill match, and target match together.
\end{itemize}

The tuned path improves all three quality-sensitive proxy dimensions while keeping format compliance fixed at 1.000. This matters because the gain is not achieved by relaxing structure; it reflects better structured decision content under the same output contract.

These proxy metrics are not a substitute for full environment-level success rates. However, for the present Assignment 3 scope they remain meaningful because they quantify structured planner-output quality under strict, directly comparable conditions: same schema, same parsing/validation contract, and same scoring logic across base and tuned paths.

\subsection{What Improved vs.\ What Did Not Improve in the Main Table}
What improved:
\begin{itemize}[leftmargin=*]
  \item Tool-use accuracy improves from 0.082 to 0.219.
  \item Target-match rate improves from 0.329 to 0.479.
  \item Strict task proxy improves from 0.068 to 0.205.
\end{itemize}

What did not change:
\begin{itemize}[leftmargin=*]
  \item Format compliance remains 1.000 in both tracks.
\end{itemize}

Hence, the quantitative result should be interpreted as a quality improvement in structured planner outputs, not as a claim of official environment-level benchmark gains.

\subsection{Tiny 3-Case Qualitative Comparison}
Table~\ref{tab:tiny-qual} reports fixed tiny side-by-side runs under matched settings.
\begin{table}[t]
\caption{Tiny 3-case diagnostic comparison under identical settings (\texttt{backend=calvin\_debug\_real}, \texttt{batch=grouped\_sequence}, \texttt{episodes=3}, \texttt{seed=42}, \texttt{verifier\_mode=verifier\_plus\_replan}, \texttt{replanner\_mode=hybrid}).}
\label{tab:tiny-qual}
\centering
\begin{tabular}{lcccc}
\toprule
Track & Parse success & Validated revisions & Accepted revised plans & Semantic rejection \\
\midrule
Stable baseline & 3/3 & 3/3 & 0/3 & 3/3 \\
Qwen2.5-VL-3B rerun & 3/3 & 3/3 & 3/3 & 0/3 \\
Qwen2.5-VL-7B rerun & 3/3 & 3/3 & 2/3 & 1/3 \\
\bottomrule
\end{tabular}
\end{table}


Required summary (exact):
\begin{itemize}[leftmargin=*]
  \item stable baseline accepted revised plans = \textbf{0/3}
  \item VL-3B = \textbf{3/3}
  \item VL-7B = \textbf{2/3}
\end{itemize}

This table is secondary qualitative evidence rather than primary quantitative evidence. Its value is controlled diagnosis: all three tracks run under matched settings, so differences can be attributed to revised-plan semantic acceptance behavior rather than setup drift.

Within this controlled setting, all tracks remain parse-stable and validation-stable (3/3), so the key separation is semantic acceptance rather than JSON formatting. VL-3B is currently the strongest secondary qualitative comparison track in this fixed set. VL-7B is runnable but does not outperform VL-3B on this tiny comparison.

The limitation is scale: a fixed 3-case set is not statistically significant and should not be interpreted as benchmark ranking. The strength is diagnostic resolution: it reveals concrete grounding and rejection patterns (\texttt{target\_absent\_from\_scene\_memory}, \texttt{drawer\_goal\_target\_mismatch}) under identical conditions.

\subsection{Short Case-Study Subsection}
\begin{table*}[t]
\caption{Side-by-side representative cases from the fixed tiny 3-case comparison. All three tracks are parse-valid and schema-validated; differences mainly appear at semantic acceptance.}
\label{tab:case-studies}
\centering
\begin{tabular}{p{0.17\textwidth}p{0.12\textwidth}p{0.20\textwidth}p{0.21\textwidth}p{0.19\textwidth}}
\toprule
Case & Episode ID & Revised-plan evidence (shortened) & Acceptance / status & One-sentence takeaway \\
\midrule
Baseline rejection case &
\texttt{episode\_0358502.npz} &
\texttt{subgoal: place\_...\_drawer; target: table; skill: place} &
parse=3/3, validated=3/3, rejected (\texttt{drawer\_goal\_target\_mismatch}), terminal label dominated by \texttt{repeated\_no\_effect\_fallback\_exhausted} &
The baseline failure is not format collapse; it is a schema-valid but semantically inconsistent revision. \\
\midrule
VL-3B accepted case &
\texttt{episode\_0358522.npz} &
\texttt{subgoal: grasp...; target: red\_block; skill: grasp} &
parse=3/3, validated=3/3, accepted (3/3 in run), terminal label still dominated by \texttt{repeated\_no\_effect\_fallback\_exhausted} &
VL-3B is strongest on this tiny set because its revised plans consistently pass semantic acceptance. \\
\midrule
VL-7B partial degradation case &
\texttt{episode\_0358502.npz} &
\texttt{subgoal: grasp\_blue\_...; target: blue\_block; skill: grasp} &
parse=3/3, validated=3/3, partially accepted (2/3), one semantic rejection (\texttt{target\_absent\_from\_scene\_memory}), terminal label dominated by \texttt{repeated\_no\_effect\_fallback\_exhausted} &
VL-7B runs successfully but does not outperform VL-3B in this controlled tiny comparison. \\
\bottomrule
\end{tabular}
\end{table*}


\textbf{Case 1: baseline bad-plan semantic rejection.}
The representative baseline episode (\texttt{episode\_0358502.npz}) shows a revised plan whose text indicates drawer intent but grounds the target to \texttt{table}. Parse and schema validation succeed, yet semantic acceptance rejects the revision with \texttt{drawer\_goal\_target\_mismatch}. This is evidence of semantic inconsistency rather than output-format failure.

\textbf{Case 2: VL-3B accepted revision.}
The representative VL-3B episode (\texttt{episode\_0358522.npz}) produces a grounded revision accepted by the semantic gate. At run level, VL-3B reaches 3/3 accepted revised plans under the same fixed setup. The remaining terminal issue is still dominated by repeated no-effect behavior, indicating downstream execution pressure after acceptance.

\textbf{Case 3: VL-7B partial degradation.}
The representative VL-7B episode (\texttt{episode\_0358502.npz}) includes one rejected revision due to \texttt{target\_absent\_from\_scene\_memory}. Run-level acceptance is 2/3, so VL-7B is operational but not stronger than VL-3B in this tiny controlled comparison.

All case descriptions are repository-backed from \texttt{docs/failure\_cases/hybrid\_replanner\_cases.md} and run metrics in the three fixed tiny run directories.

\section{Analysis and Discussion}
\subsection{Why the Tuned Path Helps}
The tuned path improves structured planner-output quality on the primary base-vs-tuned table, especially in tool-use accuracy, target-match rate, and the strict conjunction proxy. These gains suggest that minimal LoRA post-training improves action-target consistency under the existing schema constraints.

The unchanged format compliance at 1.000 is also important: it shows that the improvements are not explained by easier formatting behavior or loosened validation thresholds. Instead, the change appears in content quality dimensions while structural discipline is preserved.

\subsection{How to Interpret Tiny Qualitative Evidence}
The tiny 3-case comparison is useful because it isolates semantic acceptance behavior under one controlled setup and fixed run protocol. In this setting, all tracks are equally stable at parse/validate (3/3), so differences emerge at semantic acceptance:
\begin{itemize}[leftmargin=*]
  \item baseline: 0/3 accepted revised plans,
  \item VL-3B: 3/3 accepted revised plans,
  \item VL-7B: 2/3 accepted revised plans.
\end{itemize}

This pattern explains why VL-3B is currently the strongest secondary qualitative track and why VL-7B should be reported as runnable but not superior on this tiny set.

At the same time, tiny qualitative evidence remains diagnostic rather than statistically conclusive. The report therefore uses it to explain behavioral differences, not to claim broad generalization.

In other words, Table~\ref{tab:tiny-qual} should be read as a controlled stress test of semantic grounding quality, not as a scaled performance benchmark.

\subsection{Current Bottleneck and System Maturity}
The current failure profile is increasingly concentrated in execution-side no-effect outcomes, especially the terminal label \texttt{repeated\_no\_effect\_fallback\_exhausted}. This indicates a maturity shift: earlier failure pressure from semantically poor revisions is reduced in stronger tracks, while remaining difficulty is in turning semantically acceptable plans into effective state change.

From a course-assignment perspective, this is a meaningful intermediate maturity stage: structured output quality and semantic gating are functioning, but robust execution effectiveness remains an open engineering and evaluation challenge.

\section{Limitations}
\begin{itemize}[leftmargin=*]
  \item The tiny 3-case comparison is diagnostic only, not statistically significant benchmarking.
  \item VL-3B is a secondary qualitative track, not a replacement of the stable baseline narrative.
  \item VL-7B is not claimed to be better than VL-3B.
  \item No official CALVIN or RLBench leaderboard claim is made.
  \item Remaining failures are still execution-side/no-effect heavy.
  \item RLBench full simulator pipeline is currently blocked in this machine snapshot; fixture bridge evidence is used for cognition-layer smoke only.
\end{itemize}

\section{Conclusion}
This Assignment 3 study demonstrates an implemented agent-oriented post-training setting with structured replanning, explicit verifier feedback, and auditable recovery behavior. The strongest quantitative evidence is the base-vs-tuned proxy improvement under matched schema constraints. The tiny fixed qualitative comparison adds controlled diagnostic evidence: baseline 0/3 accepted revisions, VL-3B 3/3, VL-7B 2/3. Current scope remains conservative: useful progress in structured grounding and recovery diagnostics, with substantial remaining execution-side no-effect limitations.

\section*{Repository Evidence Paths (Appendix-Style)}
\small
This report is grounded in repository artifacts rather than an external bibliography. Load-bearing evidence paths are grouped as:
\begin{itemize}[leftmargin=*]
  \item \textbf{Main quantitative evidence:} \texttt{results/eval/planner\_base\_vs\_tuned/metrics.json}; \texttt{docs/final\_report\_assets/tables/current\_main\_results\_table.md}.
  \item \textbf{Tiny qualitative and cases:} \texttt{docs/final\_report\_assets/case\_studies/tiny\_3case\_qualitative\_summary.md}; \texttt{docs/failure\_cases/hybrid\_replanner\_cases.md}.
  \item \textbf{Project snapshot and status:} \texttt{results/reports/latest\_project\_report.json}; \texttt{results/reports/latest\_project\_dashboard.md}.
\end{itemize}

\end{document}
